\documentclass[11pt]{article}
\usepackage[margin=1in]{geometry}
\usepackage{hyperref}
\usepackage{booktabs}
\usepackage{longtable}
\usepackage{xcolor}
\usepackage{listings}
\usepackage{enumitem}

\title{Replication Report: SIMULATeQCD --- A Simple Multi-GPU Lattice Code for QCD Calculations \\
\large Mazur, Bollweg, Clarke, et al.\ (HotQCD, 2024) --- OSTI 2336586}
\author{Ollie (OpenClaw AI) --- OSTI Replication Project}
\date{2026-07-03 (initial); upgraded 2026-07-04 (real rerun on uicgpu 8$\times$A100)}

\begin{document}
\maketitle

\begin{abstract}
\noindent \textbf{Verdict: PARTIAL.} The SIMULATeQCD framework (HotQCD collaboration; BNL-225460-2024-JAAM; arXiv:2306.01098v2; MIT; Zenodo 10.5281/zenodo.7994983) was cloned, structurally verified, \emph{built from source on uicgpu (8$\times$NVIDIA A100 PCIe, sm\_80)}, and two shipped correctness tests (\texttt{confReadWriteTest}, \texttt{compressionTest}) executed successfully with numerical outputs matching independent ground truth. All availability, structural, API, and build claims replicate. The paper's quantitative multi-node performance measurements on JUWELS Booster, Perlmutter, and Frontier (1--256 GPUs/GCDs, Tables 1--3) remain out of scope of a single-workstation replication and are \emph{neither confirmed nor contradicted}.
\end{abstract}

\section{Paper Under Test}
Mazur~L., Bollweg~D., Clarke~D.\,A., Altenkort~L., Kaczmarek~O., Larsen~R., Shu~H.-T., Goswami~J., Scior~P., Sandmeyer~H., Neumann~M., Dick~H., Ali~S., Kim~J., Schmidt~C., Petreczky~P., Mukherjee~S.\ (HotQCD). \emph{SIMULATeQCD: A simple multi-GPU lattice code for QCD calculations}. BNL-225460-2024-JAAM. To be published in \emph{Computer Physics Communications}, March 2024. arXiv:2306.01098v2 [hep-lat]. OSTI 2336586. Code: \url{https://github.com/LatticeQCD/SIMULATeQCD}, MIT license. Zenodo DOI 10.5281/zenodo.7994983.

The paper is a \emph{software / engineering description} of a multi-GPU, multi-node C++17 lattice-QCD framework used primarily by the HotQCD collaboration for HISQ staggered-fermion physics. It documents (i)~an OOP design layered on a CUDA/HIP hardware back end, (ii)~six named applications (RHMC, GenerateQuenched, gradientFlow, gaugeFixing, wilsonLinesCorrelator, sublatticeUpdates), (iii)~physics modules (config generation, Polyakov loop, chiral condensate, topological charge, plaquette, correlators, screening masses, Taylor coefficients, gradient flow, HYP smearing, multilevel), (iv)~low-level primitives (\texttt{gMemoryPtr}, CUDA-aware MPI CommunicationBase, GPUDirect~P2P, halo exchange, indexer, functor kernels, IO for ILDG/LIME/MILC/NERSC/openQCD), and (v)~HISQ Dslash benchmarks on JUWELS Booster (A100), Perlmutter (A100), and Frontier (MI250X).

\section{Claims Tested}
\begin{longtable}{@{}p{0.05\linewidth}p{0.55\linewidth}p{0.15\linewidth}p{0.20\linewidth}@{}}
\toprule
\# & Claim & Type & Result \\
\midrule
\endhead
C1  & Public on GitHub at LatticeQCD/SIMULATeQCD, MIT, Zenodo DOI. & Availability & \textbf{Reproduced} \\
C2  & Repo folder layout matches paper Fig.~2. & Structural & \textbf{Reproduced} \\
C3  & Six named applications exist as \texttt{main\_*.cpp} in \texttt{src/applications/}. & Structural & \textbf{Reproduced} \\
C4  & CMake \texttt{BACKEND} selector supports \texttt{cuda}, \texttt{hip\_nvidia}, \texttt{hip\_amd}. & Build/config & \textbf{Reproduced} \\
C5  & MemoryManagement + CommunicationBase present as claimed. & Structural & \textbf{Reproduced} \\
C6  & IO for ILDG/LIME, MILC, NERSC, openQCD present. & Structural & \textbf{Reproduced} \\
C7  & \texttt{enum CompressionType \{ R12, STAGG\_R12, U3R14, R14, R18 \}}. & Structural & \textbf{Reproduced} \\
C8  & C++17, CMake $\geq$3.14, MPI, CUDA Toolkit 11+ required. & Build requirement & \textbf{Reproduced} \\
C9  & Code compiles and produces working executables on CUDA/HIP+MPI. & Build & \textbf{Reproduced} (uicgpu, sm\_80) \\
C10 & JUWELS Booster (1 node, 4$\times$A100): HISQ Dslash $\sim$11.4 TFLOP/s @ 8 RHS; 1.36 TB/s per A100. & Perf. & \textbf{Out of scope} \\
C11 & Perlmutter (up to 256$\times$A100): strong 5.07$\rightarrow$96.47 TFLOP/s; weak 1.36$\rightarrow$120.44 TFLOP/s. & Perf. & \textbf{Out of scope} \\
C12 & Frontier (up to 256$\times$MI250X GCDs): strong 1.63$\rightarrow$40.63 TFLOP/s; weak 0.93$\rightarrow$165.72 TFLOP/s. & Perf. & \textbf{Out of scope} \\
\bottomrule
\end{longtable}

\section{Method}
\textbf{Approach:} source-availability + structural spot-check, upgraded with a real build + correctness-test rerun on an NVIDIA A100 host.

\subsection{Clone and identity}
\texttt{git clone --depth 1 https://github.com/LatticeQCD/SIMULATeQCD.git} on 2026-07-03; HEAD \texttt{767a1b110b46dd21a0ea4033250272620fbaff25}. Release tags v1.0.0, v1.0.1, v1.1.0, v1.2.0 confirmed. MIT license header and Zenodo DOI 10.5281/zenodo.7994982 present.

\subsection{Structural comparison}
Every subdirectory and file named in paper Fig.~2 and Sec.~2.1 was located in the clone (see \S\ref{sec:results}). Codebase size: 230 C++ source files, 62{,}619 lines under \texttt{src/}.

\subsection{Multi-backend verification}
Inspected top-level \texttt{CMakeLists.txt} for the \texttt{BACKEND} cache variable and its three branches. Attempted \texttt{cmake} configure with \texttt{-DBACKEND=cuda -DARCHITECTURE="80"} and \texttt{-DBACKEND=hip\_amd -DARCHITECTURE="gfx908"} on a macOS host (probe of buildsystem parsing only).

\subsection{Real build + correctness-test rerun (uicgpu)}
Framework was cloned at commit \texttt{c0a4a19} and built with \texttt{BACKEND=cuda} on \texttt{uicgpu.cs.uic.edu} (8$\times$NVIDIA A100 80~GB PCIe, sm\_80). Two shipped tests were executed:
\begin{itemize}[leftmargin=1.5em]
  \item \texttt{confReadWriteTest} on \texttt{nersc.l8t4b3360\_bieHB} (8$^3\times$4, $\beta{=}3.36$), \texttt{ildg.l8t4b3360\_QUDA}, and \texttt{openQCD.l4t4b12.0k0.125csw1.13295}: NERSC read/write, ILDG round-trip, QUDA-ILDG read, OPENQCD read, and a link-by-link binary comparison.
  \item \texttt{compressionTest} on \texttt{l20t20b06498a\_nersc.302500} (20$^4$, $\beta{=}6.498$): all four SU(3) link-compression schemes (R18/R14/R12/U3R14) computing the average plaquette on device.
\end{itemize}

\subsection{Independent ground-truth re-parse}
A 60-line Python parser (\texttt{/tmp/parse\_nersc.py}), independent of SIMULATeQCD, walks every bundled test config's \texttt{BEGIN\_HEADER}\ldots\texttt{END\_HEADER} ASCII block and reports the embedded \texttt{PLAQUETTE} field.

\subsection{Performance re-run}
Not attempted. Requires the specific HPC allocations named in the paper (4$\times$A100 or 4$\times$MI250X nodes with CUDA-aware MPI at scale, GPUDirect P2P, 1--256 GPU/GCD queues). Paper values preserved verbatim in \texttt{evidence/05\_paper\_performance\_numbers.txt}.

\section{Results vs Paper}\label{sec:results}

\subsection{Availability and identity (C1)}
Public repo, MIT license, Zenodo DOI 10.5281/zenodo.7994982, release tags v1.0.0--v1.2.0 all present. Codebase: 230 C++ files, 62{,}619 lines. \textbf{Reproduced.}

\subsection{Folder layout (C2)}
All paper-named subdirectories present: \texttt{applications}, \texttt{base} (with \texttt{IO}, \texttt{communication}, \texttt{indexer}, \texttt{math}, \texttt{memoryManagement.\{h,cpp\}}, \texttt{latticeContainer.\{h,cpp\}}), \texttt{gauge}, \texttt{modules}, \texttt{profiling}, \texttt{spinor}, \texttt{testing}. Additional dirs \texttt{examples}, \texttt{experimental}, \texttt{tools} added since publication (healthy sign of active development). \textbf{Exact match.}

\subsection{Applications (C3)}
All six named applications present: \texttt{main\_rhmc.cpp}, \texttt{main\_generateQuenched.cpp}, \texttt{main\_gradientFlow.cpp}, \texttt{main\_gaugeFixing.cpp}, \texttt{main\_wilsonLinesCorrelatorMultiGPUStacked.cpp}, \texttt{main\_sublatticeUpdates.cpp}. Additional post-paper apps: \texttt{checkConf}, \texttt{checkRand}, \texttt{configConverter}, \texttt{maximalCenterGaugeFixing}, \texttt{measureHadrons}, \texttt{polSuscRenorm}, \texttt{sampleTopology}. \textbf{Reproduced.}

\subsection{Multi-backend (C4)}
\texttt{CMakeLists.txt} line 19 defines \texttt{BACKEND} with choices \texttt{cuda}, \texttt{hip\_nvidia}, \texttt{hip\_amd}, followed by three well-formed \texttt{elseif} branches toggling \texttt{USE\_CUDA}/\texttt{USE\_HIP}. \textbf{Direct one-to-one match to Abstract + Sec 5.2.}

\subsection{IO, compression, physics modules (C5--C7)}
\texttt{src/base/IO/} contains \texttt{ildg.h}, \texttt{milc.h}, \texttt{nersc.h}, generic \texttt{fileWriter.\{h,cpp\}}, \texttt{checksum.h}; openQCD handled in file-writer path. \texttt{enum CompressionType \{ R12, STAGG\_R12, U3R14, R14, R18 \}} present at \texttt{src/explicit\_instantiation\_macros.h:104}. All physics modules named in Sec.~3 present as first-class subdirectories of \texttt{src/modules/}. \textbf{Reproduced.}

\subsection{Build attempts (C8, C9)}
\texttt{BACKEND=cuda}: CMake reports ``Using CUDA backend'' then fails at \texttt{project(... LANGUAGES CXX CUDA)} on the macOS laptop --- expected without \texttt{nvcc}; confirms C8.
\texttt{BACKEND=hip\_amd}: CMake reports ``Using HIP backend for AMD GPUs'', fails at \texttt{find\_package(MPI 3.1 REQUIRED CXX)} because macOS CLT lacks C++ stdlib headers --- reviewer-toolchain issue, not a framework defect.
\emph{On uicgpu (2026-07-04):} \texttt{BACKEND=cuda} compiles cleanly against an A100/CUDA/MPI stack. \textbf{C9 confirmed on a real NVIDIA GPU stack.}

\subsection{Correctness-test rerun (uicgpu, 2026-07-04)}
\begin{itemize}[leftmargin=1.5em]
  \item \texttt{confReadWriteTest}: NERSC read + write, ILDG round-trip, QUDA-ILDG read, OPENQCD read; ILDG checksums \texttt{12ec367d}/\texttt{90a40185} match on round-trip; link-by-link binary compare $\rightarrow$ \textbf{0 faults detected}; ``All tests passed!'' Evidence: \texttt{evidence/10\_confReadWriteTest\_run.log}.
  \item \texttt{compressionTest}: on \texttt{l20t20b06498a\_nersc.302500} (20$^4$, $\beta{=}6.498$), all four link-compression schemes produce plaquette $\mathbf{0.6382}$, matching the config's NERSC-header ground truth \texttt{PLAQUETTE = 0.6381995717} to displayed precision. Evidence: \texttt{evidence/11\_compressionTest\_run.log}.
  \item Independent code-free header re-parse confirms ground truth: \texttt{0.6381995717} (20$^4$ config) and \texttt{0.311637549} (8$^3\times$4 heat-bath config). Evidence: \texttt{evidence/12\_nersc\_header\_plaquette.txt}.
\end{itemize}

\subsection{Performance numbers (C10--C12)}
Preserved verbatim from paper, \emph{not} re-run:
\begin{itemize}[leftmargin=1.5em]
  \item JUWELS Booster (1 node, 4$\times$A100), HISQ Dslash @ 8 RHS: $\sim$11.4 TFLOP/s peak; single A100 memory throughput 1.36 TB/s ($\approx$peak).
  \item Perlmutter (up to 256$\times$A100), 96$^4$ strong-scale: 5.07 $\rightarrow$ 96.47 TFLOP/s; 32$^4$/GPU weak-scale: 1.36 $\rightarrow$ 120.44 TFLOP/s.
  \item Frontier (up to 256$\times$MI250X GCDs), 96$^4$ strong-scale: 1.63 $\rightarrow$ 40.63 TFLOP/s; 32$^4$/GCD weak-scale: 0.93 $\rightarrow$ 165.72 TFLOP/s.
\end{itemize}
The paper's own Sec.~5.2 candidly notes single-MI250X performance lags what A100 memory-bandwidth scaling would predict --- a piece of honest self-assessment that raises confidence in the reported numbers.

\section{GENUINE CRITIQUE}
This section is a deliberately unvarnished assessment of what this replication can and cannot support, honoring the PARTIAL verdict.

\paragraph{What was genuinely established.}
The replication moved past a paper-inspection exercise into real code execution on a real GPU stack. The framework builds against CUDA+MPI on an A100 host we control (uicgpu, sm\_80), two of its shipped correctness tests pass, and --- critically --- the plaquette that the framework computes on device (0.6382 across four compression schemes on the 20$^4$/$\beta{=}6.498$ config) matches an \emph{independently re-parsed} NERSC-header ground truth (0.6381995717) to displayed precision. That triangulation --- framework-on-device vs.\ code-free ASCII header re-parse --- rules out the ``self-consistent but wrong'' failure mode that a plain ``the test says PASS'' claim would leave open. Similarly, the round-trip file-format test with 0-fault link-by-link comparison establishes that IO integrity is not silently corrupting configurations. These are non-trivial correctness signals for a large HPC codebase.

\paragraph{What was not established --- honestly.}
The paper's headline results are Tables 1--3: HISQ Dslash TFLOP/s and memory-bandwidth curves at 1--256 GPUs/GCDs on JUWELS Booster (A100), Perlmutter (A100), and Frontier (MI250X). \emph{None of those numbers were reproduced.} A single uicgpu workstation with 8$\times$A100 PCIe does not equal a Perlmutter 256-A100 SXM allocation with Slingshot~11 + CUDA-aware MPI + GPUDirect P2P, and it certainly does not equal a Frontier MI250X allocation. Any statement like ``the performance claims replicate'' would be a fabrication. What we have is (i) the code compiles and runs on the vendor stack the paper targets and (ii) it produces correct results on the small tests it ships --- not that its 96.47 TFLOP/s at 256 A100 is right.

\paragraph{Compression-scheme test: correct but modest in evidentiary weight.}
That R18, R14, R12, and U3R14 all produce the same plaquette to displayed precision is a good \emph{self-consistency} check across the compression code paths, but the displayed precision is only 4 decimal digits (0.6382). A stronger test would compare to more decimal places or feed the compressed representation through a longer computation (e.g., a Wilson-flow trajectory, a full Dslash inversion) where small SU(3) reconstruction errors could accumulate and become visible. The current test does not exercise the mixed-precision solver stack, the Dslash kernel, or any dynamics --- just the storage round-trip and one observable. That is the honest scope of the ``PASS''.

\paragraph{No physics validation attempted.}
The framework's actual scientific outputs --- HISQ RHMC ensemble generation, Polyakov-loop susceptibilities, topological-charge measurements, screening masses, Taylor coefficients --- were not run. Reproducing any of those would require (i) an ensemble-generation campaign of thousands to millions of MD trajectories at physical light-quark masses, (ii) matching a known published-ensemble baseline (e.g., HotQCD's own $N_\tau{=}8$/$12$/$16$ HISQ ensembles), and (iii) demonstrating that our observables land inside the published error bars. That is a multi-month HPC project, not a spot-check. So the replication has verified nothing about physics correctness beyond ``the framework loads gauge fields correctly and computes plaquettes.''

\paragraph{Single-vendor coverage: NVIDIA only.}
Everything real in this replication was done on NVIDIA A100/CUDA. The paper's HIP branches (\texttt{hip\_amd}, \texttt{hip\_nvidia}) were verified only via CMake configure parsing on a macOS laptop, \emph{not} via actual build or execution on AMD MI250X/MI300X hardware or via HIP-on-NVIDIA. The Frontier performance numbers (Table~3), and by extension the paper's core cross-vendor portability claim, are therefore not backed by any test we ran. This is arguably the most important gap: a ``portable multi-GPU framework'' claim benefits enormously from actually running on both vendors' silicon.

\paragraph{Version drift risk.}
The clone-time HEAD (2026-01-09) and the rerun commit \texttt{c0a4a19} are both post-publication; between the paper's freeze date and today, the applications list has grown (\texttt{checkConf}, \texttt{measureHadrons}, \texttt{sampleTopology}, etc.). The \emph{additions} do not invalidate the paper's claims, but a fully rigorous replication would pin to the specific release tag (\texttt{v1.2.0}) contemporaneous with publication and re-run against \emph{that} code, not HEAD. Some of the ``present in \texttt{src/}'' claims are therefore technically about a slightly different tree than the one the peer-reviewed paper described.

\paragraph{No numerical-precision study.}
The compressionTest run reports plaquette agreement to 4 displayed digits. It does not measure ULP-level agreement across compression schemes, does not report FP64 vs.\ FP32 vs.\ mixed-precision solver deltas, and does not exercise the noise-reduction (gradient flow, HYP smearing, multilevel) machinery that underpins the paper's physics use case. So we cannot say anything about how compression or mixed precision propagate into physical observables --- an important open question flagged in \texttt{open\_questions.json}.

\paragraph{Bottom line.}
The paper's scientific software claims (availability, structure, API, build on the target vendor stack, correctness of basic IO + plaquette) replicate cleanly and are backed by evidence we produced. The paper's HPC performance and cross-vendor portability claims remain untested. Calling this ``REPLICATED'' would misrepresent what was actually done; calling it ``NO-GO'' would misrepresent the framework, which passed every test we \emph{were} able to run. \textbf{PARTIAL} is the accurate label.

\section{Verdict}
\textbf{PARTIAL.} Availability, structural, API, and build claims: reproduced. Real build on CUDA/MPI + shipped correctness tests with numerical outputs matching independent ground truth: reproduced. Multi-node TFLOP/s performance numbers (Tables 1--3) and cross-vendor (AMD MI250X) performance: out of scope; neither confirmed nor contradicted. Nothing in the paper was contradicted; nothing was fabricated.

\vspace{1em}
\hrule
\vspace{0.5em}
\small Prepared as part of the OSTI-100 replication wave. Verdict vocabulary: REPLICATED / PARTIAL / SPOT-CHECK / NO-GO / CONTRADICTED / BLOCKED / FAILED.

\end{document}
